\chapter{Grandeurs, aires et volumes}
\label{corriges:17}
\section*{Exercice 17.1 -- Conversions}
$0{,}45\times10000=4500\unite{cm^2}$.
$3250\unite{cm^3}=3250\unite{mL}=3{,}25\unite L$.
$2+18/60=2{,}3\unite h$.
$90/3{,}6=25\unite{m/s}$.
Les conversions d'aires et de volumes n'utilisent pas les mêmes facteurs que les longueurs.

\section*{Exercice 17.2 -- Aires composées}
Le trapèze a pour aire $(8+14)\times5/2=55\unite{cm^2}$ ; la formule résulte de deux copies réunies en parallélogramme. La couronne est le grand disque privé du petit : $\pi7^2-\pi4^2=33\pi\unite{cm^2}$. On soustrait des aires, pas les rayons avant élévation au carré.

\section*{Exercice 17.3 -- Assemblage et débit}
Le cylindre a un volume $72\pi\unite{cm^3}$, la demi-boule un volume $\frac12\times\frac43\pi3^3=18\pi\unite{cm^3}$. Le total vaut $90\pi\unite{cm^3}$.
Le volume à fournir est $0{,}75\times1200=900\unite L$, donc la durée est $900/25=36$ minutes, sous l'hypothèse de débit constant et sans fuite.

\section*{Exercice 17.4 -- Grandeurs composées}
Le premier trajet dure $30/60=0{,}5\unite h$, le second $1\unite h$. La moyenne est $60/1{,}5=40\unite{km/h}$, et non $(60+30)/2=45$. La masse du liquide vaut $\rho V=0{,}8\times5=4\unite{kg}$. Le repère « un litre d'eau environ un kilogramme » ne se transpose pas à tous les liquides.
